% !TEX program = lualatex
% An essay for ihsan.cc/writing, written in the library. Each essay starts as a
% copy of this template (docs/publish-essay.md in the website repo). Replace the
% placeholder text, keep the macros, and say "publish essay <number> rev <R>":
% a session turns that revision into the site's markdown and opens a PR.
%
% What the site can show, and the macro for each:
%   \essayslug{...} \essaydate{YYYY-MM-DD} \essaysummary{...} \essayblurb{...}
%   \essaydraft{true|false} \essaytitle{...} \essaystandfirst{...}
%                                   the essay's details, set once at the top
%   \section{...} \subsection{...}  headings (a subsection is the smaller one)
%   \emph{} \textbf{} \texttt{} \href{url}{text}   inside a paragraph
%   itemize and enumerate           lists
%   \includesvg[caption]{figures/<name>}
%                                   a diagram: write figures/<name>.json for
%                                   diagram-maker; the build renders the PDF shown
%                                   here and the SVG the site shows
%   \essayimage{figures/<file>.png}{caption}   a photo or screenshot
%   \essayquote{...}                the one pull quote
%   \essaycode{<language>} then a Verbatim block   code
%   \essaynote{...}                 a footnote; the notes print at the end
%   \begin{essayreading} \readingitem{title}{library link}{what it covers}{pages}
%                                   further reading, library PDFs
% Anything else (a table, maths, a \footnote) has no place on the site.
\documentclass[11pt]{article}
\usepackage{housestyle}
\renewcommand{\sectionmark}[1]{\markboth{#1}{}}
\hssection{\leftmark}
\hsjustified
\setcounter{secnumdepth}{0}

\makeatletter
\newcommand{\essayslug}[1]{\def\essay@slug{#1}}
\newcommand{\essaydate}[1]{\def\essay@date{#1}}
\newcommand{\essaysummary}[1]{\def\essay@summary{#1}}
\newcommand{\essayblurb}[1]{\def\essay@blurb{#1}}
\newcommand{\essaydraft}[1]{\def\essay@draft{#1}}
\newcommand{\essaytitle}[1]{\def\essay@title{#1}\hsslug{#1}}
\newcommand{\essaystandfirst}[1]{\def\essay@standfirst{#1}}
\essayblurb{}\essaydraft{true}\def\essay@false{false}

% The title block, then what the site reads beside the text.
\newcommand{\essayhead}{%
  \hstitleblock{Essay \textperiodcentered\ \essay@date}{\essay@title}{\essay@standfirst}%
  \begin{hscallout}{For the site}
  \texttt{ihsan.cc/writing/\essay@slug}\par
  \essay@summary\par
  {\color{inkfiftyfive}In the list: \ifx\essay@blurb\@empty the summary.\else\essay@blurb\fi\ 
  \ifx\essay@draft\essay@false On /writing once published.\else
  A draft: on the preview page only.\fi}
  \end{hscallout}}

% A figure's number, then its caption under the picture. \includesvg is the
% diagram macro here (no svg package is loaded): its name is what makes the
% library keep figures/<name>.svg with the source, and the page shows the PDF.
\newcounter{essayfig}
\newcommand{\includesvg}[2][]{\stepcounter{essayfig}%
  \begin{hsfigure}{Figure \theessayfig}{#1}\hsdiagram{#2}\end{hsfigure}}
\newcommand{\essayimage}[2]{\stepcounter{essayfig}%
  \begin{hsfigure}{Figure \theessayfig}{#2}\includegraphics[width=\linewidth]{#1}\end{hsfigure}}

\newcommand{\essayquote}[1]{\hsclaim{#1}{Pull quote}}
\newcommand{\essaycode}[1]{\hslisting{Code \textperiodcentered\ #1}}

% Footnotes print together at the end, as on the site.
\newcounter{essaynote}
\def\essay@notes{}
\newcommand{\essaynote}[1]{\stepcounter{essaynote}\textsuperscript{\theessaynote}%
  \g@addto@macro\essay@notes{\item #1}}
\newcommand{\essaynotes}{\ifx\essay@notes\@empty\else
  \section*{Notes}\begin{enumerate}[leftmargin=18pt]\essay@notes\end{enumerate}\fi}

\newenvironment{essayreading}{\section*{Further reading}\begin{itemize}[leftmargin=14pt]}{\end{itemize}}
\newcommand{\readingitem}[4]{\item \href{#2}{\textbf{#1}}: #3%
  \if\relax\detokenize{#4}\relax\else\ {\color{inkfiftyfive}(#4 pages)}\fi.}
\makeatother

\begin{document}

% ---- The essay's details ----------------------------------------------------
% The slug is the address, ihsan.cc/writing/<slug>: lowercase words and hyphens.
% The summary is one line, for the essay list and link previews. The blurb is
% optional: the clause after the title in the front page's Essays list.
% Leave draft true until you want it on /writing; a draft shows only at the
% preview path.
\essayslug{a-short-slug}
\essaydate{2026-10-01}
\essaysummary{One line on what the essay is about, for the list and for link previews.}
\essayblurb{what the essay is about, in a clause.}
\essaydraft{true}
\essaytitle{The essay's title}
\essaystandfirst{One or two sentences under the title that say why the essay is worth reading.}

\essayhead

The opening paragraph. Words can be \emph{emphasised}, made \textbf{strong},
set as \texttt{code}, or linked, as in \href{https://ihsan.cc}{my site}.
A footnote goes where it is cited.\essaynote{The footnote's text, with \texttt{code} if it needs it.}

\section{A heading}

A paragraph under the heading. A diagram follows, drawn by diagram-maker from
its spec in the figures folder.

\includesvg[The diagram's caption, one or two sentences that say what to look at.]{figures/diagram}

\subsection{A smaller heading}

A list:

\begin{itemize}
\item the first item;
\item the second item.
\end{itemize}

\essayquote{The one pull quote, a sentence worth reading twice.}

\essayimage{figures/photo.png}{A photo or screenshot, and a caption that says what it shows.}

\section{Another heading}

Code, with its language:

\essaycode{sh}
\begin{Verbatim}[bgcolor=codefill]
echo "a line of code"
\end{Verbatim}

A closing paragraph.\essaynote{A second footnote.}

\essaynotes

\begin{essayreading}
\readingitem{A library document}{https://library.ihsan.cc/d/006-0032-A}{what the reader finds in it}{4}
\end{essayreading}

\end{document}
